%!TEX root = ../template.tex
%%%%%%%%%%%%%%%%%%%%%%%%%%%%%%%%%%%%%%%%%%%%%%%%%%%%%%%%%%%%%%%%%%%%
%% abstract-en.tex
%% NOVA thesis document file
%%
%% Abstract in English
%%%%%%%%%%%%%%%%%%%%%%%%%%%%%%%%%%%%%%%%%%%%%%%%%%%%%%%%%%%%%%%%%%%%

\typeout{NT FILE abstract-en.tex}%

The growing popularity of drones has been accompanied by a rising number of
unintentional incursions into restricted airspace, with documented
disruptions to airports and other critical infrastructure. This thesis
proposes a non-destructive response: a swarm of drones that encircles an
intruding, non-cooperative target and forms a visible barrier, signalling
to its operator that the airspace is restricted and encouraging the target
to leave. Here, non-cooperative means that the target shares no
control objective with the pursuers; it is not modelled as an adversary.
 
Most existing encirclement methods are planar or define encirclement as
convex containment. This work addresses the fully three-dimensional case,
where agents must be distributed evenly around the target so that no
direction is left uncovered. Three architectures are developed in
increasing generality. First, a centralized approach places agents on the
two-sphere using a Fibonacci lattice and steers them with a Lyapunov-based
geometric controller, with a Mellinger controller for low-level tracking.
Second, a decentralized approach lets agents estimate the target position
by consensus over a spatial network and redistribute themselves through
Lloyd's algorithm on Voronoi cells, with convex-hull and half-plane
implementations. Third, the formation is generalized to arbitrary surface
meshes, including a velocity-weighted variant and a barrier configuration,
and trajectory generation is replaced by a Bernstein-polynomial model
predictive control (MPC) with velocity, acceleration and collision
constraints.
 
Simulations with up to 31-agents, static and moving targets, and the
mid-flight loss of several agents show that the formations converge and
reform around the target, and that inter-agent distances stay above a
safety threshold for almost the whole simulation. Comparisons of
point-distribution, path-planning and collision-avoidance methods are also
reported. To the author's knowledge, this specific combination of
surface-based formation generation, decentralized geometric redistribution
and trajectory-level collision avoidance for a moving target has not been
systematically studied in the reviewed literature.

% Abstract keywords
\keywords{
Multi-UAV encirclement \and
Formation control \and
Lyapunov-based controller \and
Unmanned Aerial Vehicle\and
Unmanned Ground Vehicle \and
Model Predictive Control \and
Collision avoidance \and
Voronoi diagrams and Lloyd's algorithm \and
Path planning
}

